最优控制那一章里，动力学和代价都写在纸上，剩下的是求解。这里把纸拿走。智能体只能一遍遍地观察状态、选一个动作、收到一个数、落到下一个状态，再从这些局部交互里凑出一套长期策略。麻烦不止是``公式不知道''：当前这个动作决定了下一批数据长什么样，学习者同时是统计者和数据生成过程的一部分，统计学里那条``数据与参数无关''的默认前提从一开始就不成立。

整个单元可以读成一串放松。每放松一个条件，能处理的问题多一类，能兑现的保证少一层。

\begin{center}
\begin{tikzpicture}[>=stealth,every node/.style={font=\small}]
\node[draw,align=center] (a) at (0,4.5) {动态规划\\模型已知};
\node[draw,align=center] (b) at (0,3.0) {采样与自举\\模型未知};
\node[draw,align=center] (c) at (0,1.5) {函数逼近\\状态空间太大};
\node[draw,align=center] (d) at (0,0) {离策略\\数据来自别人};
\draw[->] (a) -- (b);
\draw[->] (b) -- (c);
\draw[->] (c) -- (d);
\node[right,align=left] at (1.6,3.75) {换来：不必写出转移矩阵\\交出：只在访问过的地方有效};
\node[right,align=left] at (1.6,2.25) {换来：连续与高维状态可表示\\交出：逐状态的独立表示，以及压缩性};
\node[right,align=left] at (1.6,0.75) {换来：历史数据与他人数据可复用\\交出：更新分布与评估分布的对齐};
\end{tikzpicture}
\end{center}

\emph{一级一级往下放松，右边那两行是每一级的收支；读这一章时最该记住的是自己站在第几级。}

前四篇走完第一级。\hyperref[sec:16-Control-and-Reinforcement-Learning/02-Reinforcement-Learning/01]{``MDP 把决策问题写成随机过程''} 把交互写成状态、动作、转移、奖励与折扣率，重点在 Markov 性对状态提出了什么要求——状态里少放一样东西，后面所有推导的前提就塌了。\hyperref[sec:16-Control-and-Reinforcement-Learning/02-Reinforcement-Learning/02]{``Bellman 方程与值迭代、策略迭代''} 从``最优策略的尾部仍是最优''推出递归，用压缩映射解释值迭代为何一定收敛，这是全章唯一一个不打折扣的保证，后面的每一步都在消耗它。\hyperref[sec:16-Control-and-Reinforcement-Learning/02-Reinforcement-Learning/03]{``Monte Carlo 与随机逼近把回报变成增量更新''} 拿走模型：把回报看成价值的随机样本，用 Robbins--Monro 式的步长条件把样本平均变成增量更新。\hyperref[sec:16-Control-and-Reinforcement-Learning/02-Reinforcement-Learning/04]{``从动态规划到 Q 学习''} 用一次真实转移代替 Bellman 期望，比较 Sarsa、Expected Sarsa 与表格 Q 学习，并说明它们的收敛条件为什么写在``表格''这两个字上。

第五到第八篇处理第二级和第三级，也是这一章最容易踩坑的一段。\hyperref[sec:16-Control-and-Reinforcement-Learning/02-Reinforcement-Learning/05]{``线性价值函数逼近求解投影 Bellman 方程''} 说明参数化之后被求解的已经是另一个方程，并区分价值误差、Bellman 误差与投影不动点这三个容易混为一谈的量。\hyperref[sec:16-Control-and-Reinforcement-Learning/02-Reinforcement-Learning/06]{``Deadly Triad 让自举误差在分布外放大''} 是这一章的枢纽：函数逼近、自举、离策略三者任取其二都有收敛结论，三者齐全时期望更新本身可能没有吸引不动点，两个状态的例子就足以把它演出来。\hyperref[sec:16-Control-and-Reinforcement-Learning/02-Reinforcement-Learning/07]{``DQN 用慢目标与回放缓和训练耦合''} 把目标网络和经验回放放回这个三角里看：它们各自放慢一条耦合，一角也没拆掉，属于工程缓解。\hyperref[sec:16-Control-and-Reinforcement-Learning/02-Reinforcement-Learning/08]{``Off-policy 校正用权重交换分布偏差与方差''} 把``这批数据来自谁''写成概率比率，权重于是按轨迹长度指数集中，截断与自归一化都是在偏差与方差之间挑一个点。

最后两篇换一条路线。\hyperref[sec:16-Control-and-Reinforcement-Learning/02-Reinforcement-Learning/09]{``策略梯度直接优化动作分布''} 不再绕道值函数，用 log-derivative 恒等式把``改变分布''翻译成``给样本加权重''，省掉了连续动作上的 \(\max\)，换来正比于轨迹长度的方差。\hyperref[sec:16-Control-and-Reinforcement-Learning/02-Reinforcement-Learning/10]{``Actor--Critic 在偏差与方差之间搭桥''} 用学出来的价值函数替掉整条回报，多步目标、TD(\(\lambda\)) 与 GAE 是同一条偏差---方差轴上的不同取点，双时间尺度则解释了 Actor 与 Critic 为何要一快一慢。

有两处与别的章相通。策略每一刻都在主动改变动作的分配机制，进而改变未来的状态分布，对既有数据做条件化表达不了这种干预；\hyperref[ch:120]{``因果推断：预测世界不等于改变世界''} 讨论的是一次性干预下的可识别性与重叠条件，这里把同一件事推进到序贯场景，历史依赖、动态混杂和支持缺口会一并出现。若只行动一次，\hyperref[ch:092]{``统计决策与可信预测：模型给出概率之后怎样行动''} 里的后验风险与 minimax 已经够用；强化学习多出来的困难是行动会创造下一步的数据，眼前风险最小不等于长期回报最大。至于模型已知时的那一套，\hyperref[ch:135]{``最优控制：从反馈稳定到全局最优''} 给出了图上第一级的完整答案，两章对照着读，能看清``学出来的''和``解出来的''差在哪里。

\subsection{读之前先把四种失败分开}

这一章的算法出问题时，症状常常长得一样：回报曲线不动，或者数字飞了。病因至少有四类，补救动作完全不同。探索不足是覆盖问题，某些状态动作从未被访问，任何算法都变不出没有的信息。方差过大是采样问题，加大批量、缩短有效视界、减掉基线都能缓解。偏差是逼近问题，Critic 不准或截断过狠，加样本量不管用。发散是动力学问题，期望更新本身就没有吸引不动点，降学习率只能推迟。看到 Q 值飞涨就去调学习率，看到曲线不动就去加训练步数，多半是在治错的那一层。

读任何一条收敛结论，都要连同它的条件一起读：状态空间是否有限、是否折扣、访问频率是否被保证、步长是否满足 Robbins--Monro 条件、策略在评估期间是否固定、数据是 on-policy 还是 off-policy。这六条里换掉任何一条，结论都可能作废，而符号一个都不用改。十篇共享一套记号：状态 \(s\)、动作 \(a\)、奖励 \(r\)、折扣率 \(\gamma\)、状态值 \(V^\pi\)、动作值 \(Q^\pi\)；全章限定在离散时间与折扣回报上，有限时域、平均回报、部分可观测与连续时间都需要重写而不是换符号。

最短的主线是前四篇，把问题写清、写出递归、换成采样、变成增量算法。第五到第八篇适合第二遍连读，它们回答同一个问题：加了函数逼近又复用了异分布数据之后，为什么会不稳定。最后两篇可以独立进入，前提是先接受一件事——那里的方差不是实现不够好造成的，它写在估计量的定义里。
